%% linearprogram.tex
\section{The Linear Program and Its Geometry}\label{sec:lp}

The flux constraint retaining growth dilution,
Equation~\eqref{eq:dilution-constraint},
together with its steady-state special case
Equation~\eqref{eq:sv0}, is a statement of conservation: it does not by
itself pick out any particular flux
distribution among all those consistent with it. Turning conservation
into a prediction requires two further ingredients. The first is a box of
flux bounds, $\boldsymbol\ell\le\hat{\mathbf v}\le\mathbf u$, one lower and
one upper bound per reaction, encoding thermodynamic reversibility
($\ell_j<0$ permitted), enzyme capacity, and substrate availability; how
those bounds are set is taken up next. The
second is an objective, a linear functional $\mathbf c^\top\hat{\mathbf
v}$ of the flux vector that states what the cell, or the modeler, is
presumed to optimize. With both in hand the mole-balance constraint
becomes a linear program:
\begin{equation}\label{eq:fba-lp}
\begin{aligned}
  \max_{\hat{\mathbf v}\,\in\,\mathbb R^n}\quad & \mathbf c^\top\hat{\mathbf v}\\
  \text{subject to}\quad & \mathbf S\hat{\mathbf v} = \mathbf 0\ \ (\text{or }\mu\mathbf x),\\
  & \boldsymbol\ell \le \hat{\mathbf v} \le \mathbf u.
\end{aligned}
\end{equation}
This linear program is the complete textbook formulation of flux balance
analysis \cite{Orth2010}. The objective vector $\mathbf c$ is a
choice, not a derived quantity: setting $c_j=1$ at the index of a biomass
pseudo-reaction, one that drains precursors in the ratios required to
build new cell mass, casts the LP as growth-rate maximization, the
default assumption for a cell under selection pressure to divide. Setting
$c_j=1$ instead at the index of a product-export exchange reaction casts
the same LP as yield maximization for a chosen metabolite, the framing
relevant to a bioprocess engineer optimizing a production strain rather
than modeling wild-type physiology. These cases have the same constraint
set and differ only in $\mathbf c$.

Linear programming addresses the underdetermination of the constraint
set. A genome-scale reconstruction lists every reaction
annotated to a genome, and the reaction count $n$ typically exceeds the
metabolite count $m$, because
a single metabolite typically participates in several reactions while a
single reaction typically touches only a handful of metabolites. The
stoichiometric matrix $\mathbf S\in\mathbb R^{m\times n}$ therefore has
many more columns than rows, and its null space $\mathcal N(\mathbf S) =
\{\hat{\mathbf v}\in\mathbb R^n : \mathbf S\hat{\mathbf v}=\mathbf 0\}$ is
nontrivial: a whole subspace, not a point, satisfies the steady-state
constraint. Every flux vector in $\mathcal N(\mathbf S)$ conserves mass at
every internal node exactly as well as every other; conservation alone
cannot distinguish among them. The flux bounds cut this unbounded
subspace down to a bounded polytope, and the objective in
Equation~\eqref{eq:fba-lp} selects a single vertex, or a face, of that
polytope, but the underlying indeterminacy, that most of the space
allowed by stoichiometry is not resolved by stoichiometry, remains after
an optimum has been found. This progression from conservation to bounds
to optimization is summarized in Figure~\ref{fig:lp-geometry}.

\begin{figure}[htbp]
  \centering
  \includegraphics[width=\linewidth]{lp_geometry.pdf}
  \caption{Geometry of flux balance analysis in a two-flux projection.
    (a) Conservation restricts steady-state fluxes to the null space
    $\mathcal N(\mathbf S)$; the rays represent admissible steady-state
    directions in the displayed projection. (b) Lower and upper flux bounds
    intersect that space to form the convex feasible polytope $\mathcal P$.
    (c) Parallel level sets of the linear objective
    $\mathbf c^{\mathsf T}\hat{\mathbf v}$ move in the objective direction
    until they support the polytope at an optimum. That optimum may be a
    unique vertex or an entire face. In the latter case, fluxes can vary while
    retaining the same objective value, and flux variability analysis projects
    the optimal face onto a reaction-specific interval, illustrated here for
    $\hat v_1$.}
  \label{fig:lp-geometry}
\end{figure}

The singular value decomposition makes the geometry of that
indeterminacy explicit. Writing $\mathbf S = \mathbf U\boldsymbol\Sigma
\mathbf V^\top$, with $\mathbf U\in\mathbb R^{m\times m}$ and $\mathbf
V\in\mathbb R^{n\times n}$ orthogonal and $\boldsymbol\Sigma\in\mathbb
R^{m\times n}$ diagonal with nonnegative entries, the columns of $\mathbf
V$ associated with zero singular values span $\mathcal N(\mathbf S)$, the
right null space: every feasible steady-state flux mode is a linear
combination of these columns, and reading Equation~\eqref{eq:fba-lp}
through this lens, the LP is a search over coefficients of that
combination rather than over the raw flux vector. The columns of $\mathbf
U$ associated with zero singular values span the left null space of
$\mathbf S$, equivalently the null space of $\mathbf S^\top$; a vector
$\mathbf y$ in this space defines a conserved moiety, a linear combination
$\mathbf y^\top\mathbf x$ of metabolite concentrations, such as a bound
cofactor pair sharing a fixed pool total, that is invariant under the
reaction network regardless of which flux vector in $\mathcal N(\mathbf
S)$ is realized. For a genome-scale reconstruction the rank of $\mathbf S$
is smaller than the reaction count, $\mathrm{rank}(\mathbf S)\ll n$,
so the right null space, and with it the space of flux distributions the
constraints alone cannot distinguish, is large \cite{Bordbar2014}.
That indeterminacy can be read reaction by reaction. Because an LP
optimum need not be unique, and because many reactions can
vary freely without changing the objective value at all, flux variability
analysis reports a range rather than a point for each flux: with the
objective fixed at its optimal value $Z^\star$ (or at a specified
fraction of it), each $v_j$ is in turn minimized and maximized subject to
Equation~\eqref{eq:fba-lp}'s constraints \cite{Heirendt2019}. A reaction
whose minimum and maximum coincide is pinned by the network topology and
bounds at the optimal objective value used for that calculation, while a
reaction with a wide range
is left essentially free; a different objective can produce a different
feasible face and a different range for the same reaction, and the
pattern of which fluxes fall into which
category is exactly the information that a single optimal flux vector,
taken alone, cannot report.
